\documentclass[aspectratio=169]{beamer}
\input{header}
\coursecode{JKSSB}

\title{Malware}
\subtitle{Lesson 42 --- The Malware Umbrella}
\author{JKSSB Computer Awareness}
\date{\today}

\begin{document}
\frame{\titlepage}

\begin{frame}{What ``Malware'' Means}
  \begin{itemize}
    \item \key{Malware} is short for \key{malicious software}: any program
      written to harm a computer, steal from it, or misuse it.
    \item The official syllabus names \key{computer virus and latest
      anti-virus} explicitly; the exam also tests the neighbours: worm,
      Trojan, ransomware, spyware, adware, and more.
    \item One umbrella word, many behaviours: some \key{damage} files, some
      \key{steal} data, some \key{demand money}, some \key{hide} the attacker.
    \item The exam rarely asks for stories; it asks \key{which behaviour
      belongs to which name}.
  \end{itemize}
  \begin{block}{The one idea}
    Learn each type by its method: how it spreads, what it wants, and what
    gives it away.
  \end{block}
\end{frame}

\begin{frame}{The Umbrella Map}
  \begin{center}
    \begin{tabular}{p{0.20\textwidth}p{0.34\textwidth}p{0.34\textwidth}}
      \toprule
      \textbf{Type} & \textbf{What it does} & \textbf{One-line sign} \\
      \midrule
      Virus & Attaches to files, needs a click & Infected attachment runs \\
      Worm & Self-spreading over networks & Travels alone, no host file \\
      Trojan & Pretends to be useful software & User installs it willingly \\
      Ransomware & Locks files, demands payment & ``Pay to unlock'' screen \\
      Spyware / keylogger & Watches and steals quietly & Stolen passwords \\
      Adware & Floods the screen with ads & Pop-ups everywhere \\
      Rootkit / backdoor & Hides the attacker, keeps entry open & Nothing visible \\
      Bot / logic bomb / cryptojacking & Misuses the machine & Slow PC, timed trigger \\
      \bottomrule
    \end{tabular}
  \end{center}
  \vspace{0.25cm}
  \muted{Details follow; this table is the map to keep in mind.}
\end{frame}

\begin{frame}{Virus}
  \begin{itemize}
    \item A \key{virus} is malware that \key{attaches itself to a host}
      program or file and \key{needs human action} to run.
    \item It spreads when the infected file is opened, shared, or copied:
      email attachment, pen drive, downloaded file.
    \item It can damage, delete, or corrupt files once it runs.
    \item Key exam point: a virus \key{cannot spread by itself}; it needs a
      host file and a user to launch it.
  \end{itemize}
  \begin{block}{Keep the pair straight}
    Virus = needs a host file \key{and} a human click to spread.
  \end{block}
\end{frame}

\begin{frame}{Worm}
  \begin{itemize}
    \item A \key{worm} is malware that \key{spreads by itself over a
      network}, with no host file and no human click.
    \item It copies itself from one machine to another through network
      connections and security gaps.
    \item It mainly harms by \key{consuming network and system resources}:
      the network slows down, machines slow down, services jam.
    \item Key exam point: a worm is \key{self-replicating and standalone};
      that is what separates it from a virus.
  \end{itemize}
  \begin{block}{Keep the pair straight}
    Worm = travels alone over the network; needs \key{no} host and
    \key{no} click.
  \end{block}
\end{frame}

\begin{frame}{Trojan (Trojan Horse)}
  \begin{itemize}
    \item A \key{Trojan} (Trojan horse) \key{pretends to be useful software}:
      a free game, a useful tool, a harmless attachment.
    \item The user \key{installs or opens it willingly}, and the hidden
      malicious part runs with it.
    \item Unlike a virus or worm, a Trojan \key{does not replicate itself};
      it relies on tricking the user once per machine.
    \item Once inside, it can steal data, open a \key{backdoor}, or install
      further malware such as spyware or ransomware.
  \end{itemize}
  \begin{alertblock}{Trap alert}
    Trojan = disguised as legitimate; it \key{does not self-replicate}.
  \end{alertblock}
\end{frame}

\begin{frame}{Virus vs Worm vs Trojan}
  \begin{center}
    \begin{tabular}{llll}
      \toprule
      \textbf{Point} & \textbf{Virus} & \textbf{Worm} & \textbf{Trojan} \\
      \midrule
      Needs host file & Yes & No & No (pretends to be one) \\
      Spreads by itself & No & Yes & No \\
      Needs user action & Yes & No & Yes (install/open) \\
      Main harm & File damage & Network slowdown & Hidden entry/theft \\
      \bottomrule
    \end{tabular}
  \end{center}
  \vspace{0.35cm}
  \begin{alertblock}{The core exam table}
    Virus needs host + click; worm needs neither; Trojan tricks the user and
    does not replicate.
  \end{alertblock}
\end{frame}

\begin{frame}{Ransomware}
  \begin{itemize}
    \item \key{Ransomware} \key{locks} the victim's files by
      \key{encryption} or locks the whole screen, then \key{demands payment}
      (ransom) for the unlock key.
    \item The signature sign is the \key{``pay to unlock'' message} with a
      countdown or payment address.
    \item It typically arrives like a Trojan: a harmful attachment or link
      the user opens.
    \item Defence the exam expects: regular \key{offline backups}, updated
      software, and \key{anti-virus} software --- paying does not guarantee
      recovery.
  \end{itemize}
  \begin{exampleblock}{Office desktop example}
    A clerk opens an invoice attachment; files turn unreadable and a message
    demands payment to restore them --- that is ransomware.
  \end{exampleblock}
\end{frame}

\begin{frame}{Spyware}
  \begin{itemize}
    \item \key{Spyware} secretly \key{watches} what the user does and sends
      it to someone else: browsing habits, logins, personal details.
    \item It \key{hides quietly} and tries not to damage anything, so the
      victim notices nothing.
    \item It often arrives bundled with free downloads or deceptive
      installers.
    \item Purpose is \key{theft by observation}: credentials, bank details,
      and personal data.
  \end{itemize}
  \begin{block}{Keep the pair straight}
    Spyware = silent watcher; its goal is stolen information, not visible
    damage.
  \end{block}
\end{frame}

\begin{frame}{Adware and Keylogger}
  \begin{itemize}
    \item \key{Adware} floods the machine with \key{unwanted advertisements}:
      pop-ups, banners, redirected searches.
    \item Its motive is usually money from ads, but heavy adware also
      \key{slows} the machine and can carry spyware with it.
    \item A \key{keylogger} records \key{every key pressed} and sends the
      log to the attacker --- a focused spying tool for passwords and card
      numbers.
    \item Keyloggers may be software, or small \key{hardware} devices fitted
      between keyboard and computer.
  \end{itemize}
  \begin{block}{Neighbour line}
    Adware announces itself with ads; spyware and keyloggers hide and steal.
  \end{block}
\end{frame}

\begin{frame}{Spyware vs Adware}
  \begin{center}
    \begin{tabular}{lll}
      \toprule
      \textbf{Point} & \textbf{Spyware} & \textbf{Adware} \\
      \midrule
      Visible to user & No (hides) & Yes (pop-ups, ads) \\
      Main goal & Steal data quietly & Show ads for money \\
      Gives itself away & Stolen accounts & Flooded screen \\
      \bottomrule
    \end{tabular}
  \end{center}
  \vspace{0.35cm}
  \begin{alertblock}{The second core table}
    Spyware hides and steals; adware shows itself and advertises.
  \end{alertblock}
\end{frame}

\begin{frame}{Rootkit and Backdoor}
  \begin{itemize}
    \item A \key{rootkit} is malware that \key{hides} itself and other
      malware deep inside the system, often with administrator-level control.
    \item Its job is \key{concealment}: files, processes, and network
      activity are kept invisible to the user and to ordinary tools.
    \item A \key{backdoor} is a \key{secret entry path} left into the
      system, bypassing normal login and checks.
    \item Trojans and rootkits commonly \key{install backdoors} so the
      attacker can return later without being noticed.
  \end{itemize}
  \begin{block}{Keep the pair straight}
    Rootkit = hides the intruder; backdoor = the hidden door the intruder
    keeps.
  \end{block}
\end{frame}

\begin{frame}{Bot, Botnet, Logic Bomb, Cryptojacking}
  \begin{itemize}
    \item A \key{bot} is a compromised machine controlled remotely; a
      \key{botnet} is a \key{network of such machines} (robot + network)
      obeying one master, used for spam and large-scale attacks.
    \item A \key{logic bomb} is malicious code that \key{waits for a trigger}:
      a date, an event, or an action --- then strikes.
    \item \key{Cryptojacking} secretly uses the victim's computer to
      \key{mine cryptocurrency}, so the machine runs hot and slow while the
      attacker earns.
    \item Common thread: the machine is \key{misused} while the owner notices
      only slowness or a timed surprise.
  \end{itemize}
\end{frame}

\begin{frame}{How Malware Reaches You --- and How to Stop It}
  \begin{itemize}
    \item Common entry routes: harmful \key{email attachments and links},
      infected pen drives, free-software bundles, and unpatched systems.
    \item \key{Anti-virus} software detects, blocks, and removes known
      malware; keep it and the operating system \key{updated}.
    \item Practical habits: do not open unknown attachments, scan pen
      drives, download from trusted sources, and keep \key{offline backups}.
    \item A \key{firewall} filters network traffic; it is a barrier, not a
      cleaner --- it does not remove malware already inside.
  \end{itemize}
  \begin{block}{Exam line}
    Anti-virus removes malware; firewall filters traffic; backups defeat
    ransomware.
  \end{block}
\end{frame}

\begin{frame}{Trap Alert: Malware Facts to Keep Straight}
  \begin{itemize}
    \item \key{Virus} needs a host + click; \key{worm} needs neither;
      \key{Trojan} disguises itself and does not replicate.
    \item \key{Spyware} hides and steals; \key{adware} shows ads.
    \item \key{Ransomware} encrypts and demands payment; \key{keylogger}
      records keystrokes.
    \item \key{Rootkit} hides; \key{backdoor} is the secret re-entry;
      \key{botnet} is a network of controlled machines.
    \item \key{Logic bomb} waits for a trigger; \key{cryptojacking} mines
      cryptocurrency on your machine.
  \end{itemize}
\end{frame}

\begin{frame}{Lesson 42: Recall}
  \begin{itemize}
    \item Malware = malicious software; the umbrella covers virus, worm,
      Trojan, ransomware, spyware, adware, rootkit, keylogger, bot/botnet,
      backdoor, logic bomb, and cryptojacking.
    \item Virus vs worm vs Trojan: host + click, neither, disguise without
      replication.
    \item Spyware vs adware: silent theft vs visible ads; keylogger is a
      focused spying tool.
    \item Rootkit hides the intruder; backdoor keeps the door open; botnet is
      a controlled network; logic bomb waits for its trigger; cryptojacking
      mines on your machine.
    \item Protection: updated anti-virus, firewall, careful clicks, trusted
      downloads, offline backups.
  \end{itemize}
\end{frame}
\end{document}
